ក្រសួងសុខាភិបាល សូមជម្រាបជូនសាធារណជនឱ្យបានជ្រាបថា ចាប់តាំងពីត្រីមាសទី៣ នៃឆ្នាំ២០២៥ រហូតមកដល់ពេលនេះ ប្រទេសកម្ពុជាបានសម្រេចសមិទ្ធផលជាប្រវត្តិសាស្ត្រ ក្នុងការលុបបំបាត់ការចម្លងជំងឺ គ្រុនចាញ់ក្នុងស្រុក តាមរយៈការអនុវត្តផែនការយុទ្ធសាស្ត្រជាតិលុបបំបាត់ជំងឺគ្រុនចាញ់ ឆ្នាំ២០១១-២០២៥ ដែលជាគំនិតផ្តួចផ្តើមប្រកបដោយគតិបណ្ឌិតរបស់ \textbf{សម្តេចអគ្គមហាសេនាបតីតេជោ ហ៊ុន សែន} អតីតនាយករដ្ឋមន្ត្រី និងត្រូវបានបន្តដោយ \textbf{សម្តេចមហាបវរធិបតី ហ៊ុន ម៉ាណែត} នាយករដ្ឋមន្ត្រី នៃ ព្រះរាជាណាចក្រកម្ពុជា។ សមិទ្ធផលនេះ បានមកពីប្រសិទ្ធភាពនៃការអនុវត្តយុទ្ធសាស្ត្រក្នុងការរុករក ការ ព្យាបាល ការបង្ការ ការតាមដាន និងការធ្វើអន្តរាគមន៍លើករណីជំងឺគ្រុនចាញ់បានឆាប់ និងទាន់ពេលវេលា ដែល ដឹកនាំអនុវត្តប្រកបដោយការទទួលខុសត្រូវខ្ពស់របស់ក្រសួងខាភិបាល រួមជាមួយអ្នកពាក់ព័ន្ធនៅគ្រប់កម្រិតនានា មានដូចជា អង្គការសុខភាពពិភពលោក (WHO) ព្រមទាំងក្រសួង ស្ថាប័ននានា អាជ្ញាធរ ដៃគូអភិវឌ្ឍន៍ វិស័យ ឯកជន និងសហគមន៍ទូទាំងប្រទេស ដើម្បីរួមចំណែកជាមួយរាជរដ្ឋាភិបាល ក្នុងការលើកកម្ពស់សុខុមាលភាព របស់ប្រជាជន និងការអភិវឌ្ឍសេដ្ឋកិច្ច-សង្គម។\par
ក្នុងបរិបទដែលប្រទេសកម្ពុជា កំពុងស្ថិតក្នុងដំណាក់កាល ទប់ស្កាត់ការកកើតជំងឺគ្រុនចាញ់ឡើងវិញ-ដើម្បីឈានទៅទទួលវិញ្ញាបនបត្រលុបបំបាត់ជំងឺគ្រុនចាញ់នាឆ្នាំ២០៣០ ពីអង្គការសុខភាពពិភពលោក។ ប្រទេស កម្ពុជាអាចនៅប្រឈមហានិភ័យ ចំពោះការនាំចូលជំងឺគ្រុនចាញ់ពីក្រៅប្រទេសផ្សេងទៀត ដែលកំពុងមានការ ចម្លងជំងឺគ្រុនចាញ់ តាមរយៈការធ្វើដំណើរផ្លាស់ទីរបស់ប្រជាជនប្រឈម ចូលមកប្រទេសកម្ពុជា រួមទាំងគ្រុន ចាញ់នាំចូលតាមរយៈសត្វស្វា។ ក្រសួងសុខាភិបាល សូមជម្រាបជូនសាធារណជនឱ្យបានជ្រាបថា ក្នុងរយៈពេល ៦ខែដើមឆ្នាំ២០២៦ ប្រទេសកម្ពុជាមិនបានរកឃើញករណីជំងឺគ្រុនចាញ់ឆ្លងថ្មី ក្នុងស្រុកទេ បានន័យថាគ្រុន ចាញ់បានថយចុះ {១០០\%} និងពុំមានករណីស្លាប់ដោយសារជំងឺគ្រុនចាញ់ចាប់តាំងពីឆ្នាំ២០១៨ រហូតមកដល់ បច្ចុប្បន្ននេះ។\par
សព្វថ្ងៃ ប្រទេសកម្ពុជាកំពុងដំណើរការប្រព្រឹត្តទៅយ៉ាងសកម្មដ៏មានប្រសិទ្ធភាពក្នុងការតាមដាន និងរុក រកករណីជំងឺគ្រុនចាញ់ តាមរយៈប្រព័ន្ធរបាយការណ៍រហ័សទាន់ពេលវេលា ២៤ម៉ោង (real-time) មាននៅគ្រប់ បណ្តាភូមិសាស្ត្រ ក្នុងសហគមន៍ជាអ្នកស្ម័គ្រចិត្តព្យាបាលជំងឺគ្រុនចាញ់ភូមិ និងបណ្តាញសេវាសុខាភិបាល សាធារណៈ ជាប៉ុស្តិ៍សុខភាព មណ្ឌលសុខភាព មន្ទីរពេទ្យបង្អែក មន្ទីរពេទ្យខេត្ត និងមន្ទីរពេទ្យជាតិ ដើម្បីឈាន ទៅការឆ្លើយតបរហ័ស ២៤ម៉ោង និងការទប់ស្កាត់ការកើតឡើងវិញនៃជំងឺគ្រុនចាញ់បានទាន់ពេលវេលា។ ជា លទ្ធផលនៃប្រសិទ្ធភាពប្រព័ន្ធព័ត៌មានជំងឺគ្រុនចាញ់នេះ ជំងឺគ្រុនចាញ់ឆ្លងក្នុងស្រុកថ្មីនាឆ្នាំ២០២៦ ក៏កំពុងត្រូវ បានរកឃើញថា មិនមានករណីជំងឺគ្រុនចាញ់សង្ស័យណាមួយលេចចេញតាមបណ្តាព្រំដែនប្រទេសកម្ពុជានោះទេ។\par
អាស្រ័យហេតុនេះ ក្រសួងសុខាភិបាលសូមលើកទឹកចិត្តដល់សាធារណជន ចូលរួមបង្ការជំងឺគ្រុនចាញ់ តាមរយៈការការពារខ្លួនកុំឱ្យមូសខាំ ត្រូវសម្រាន្តក្នុងមុងគ្រប់ពេល សូមស្វែងរកសេវាពិគ្រោះប្រឹក្សាយោបល់ ឬធ្វើតេស្តនិងព្យាបាលឱ្យបានទាន់ពេលវេលា តាមរយៈអ្នកស្ម័គ្រចិត្តភូមិព្យាបាលជំងឺគ្រុនចាញ់ ឬមូលដ្ឋាន សុខាភិបាលសាធារណៈដែលនៅជិតបំផុត ប្រសិនបើសង្ស័យថាបានឆ្លងជំងឺគ្រុនចាញ់។ ក្រសួងសុខាភិបាលនឹង បន្តធ្វើការជូនដំណឹង អំពីព័ត៌មានទាក់ទងនឹងបញ្ហាសុខភាពសាធារណៈ តាមបណ្តាញសង្គមផ្លូវការរបស់ក្រសួង សុខាភិបាល \href{https://www.moh.gov.kh}{\nolinkurl{www.moh.gov.kh}} ក៏ដូចជាទំព័រហ្វេសប៊ុកផ្លូវការរបស់នាយកដ្ឋានប្រយុទ្ធនឹងជំងឺឆ្លង និងគេហទំព័រ \href{https://www.cdcmoh.gov.kh}{\nolinkurl{www.cdcmoh.gov.kh}} និងគេហទំព័ររបស់មជ្ឈមណ្ឌលជាតិប្រយុទ្ធនឹងជំងឺគ្រុនចាញ់ ប៉ារ៉ាស៊ីតសាស្ត្រ និងបាណកសាស្ត្រ \href{https://www.com.gov.kh}{\nolinkurl{www.com.gov.kh}} ។
